\documentclass[11pt]{article}

\usepackage[margin=0.9in]{geometry}
\usepackage{amsmath,amssymb,amsthm,mathtools}
\usepackage{booktabs}
\usepackage{microtype}
\usepackage[hidelinks]{hyperref}
\usepackage{xurl}
\emergencystretch=2em

\newtheorem{theorem}{Theorem}
\newtheorem{lemma}[theorem]{Lemma}
\newtheorem{proposition}[theorem]{Proposition}
\newtheorem{corollary}[theorem]{Corollary}
\theoremstyle{remark}
\newtheorem{remark}[theorem]{Remark}

\newcommand{\R}{\mathbb R}
\newcommand{\Z}{\mathbb Z}
\newcommand{\E}{\mathbb E}
\newcommand{\OPT}{\operatorname{OPT}}
\newcommand{\ALG}{\operatorname{ALG}}
\newcommand{\CPB}{C_{\mathrm{PB}}}

\title{A Constant-Competitive Algorithm for Dynamic\\
Mixture-of-Experts Serving}
\author{Ian D'Ambrosio\\
\small Nth Research Collective\\
\small \texttt{ian@nthresearch.org}}
\date{11 August 2026}

\begin{document}
\maketitle

\begin{abstract}
Huang, Lou, and Xiao introduced Dynamic Mixture-of-Experts Serving and gave an
\mbox{$O(\sqrt{\log k})$}-competitive randomized algorithm for its integral primal
problem, where $k$ is the number of replica GPUs beyond the mandatory copy of
each expert.  Their matching lower barrier applies to an auxiliary dual and
leaves the primal order open.  We prove that the randomized primal competitive
ratio is in fact $\Theta(1)$ for arbitrary numbers of experts.  The upper bound
reduces reciprocal-max service costs to chasing positive bodies with covering
row sparsity two.  A finite tangent envelope approximates each reciprocal
epigraph within a constant factor, summable positive resets convert accumulated
service into movement, and a nonexpansive balanced projection removes the
positive-body algorithm's resource augmentation.  Combining the resulting
fractional path with Lazy Threshold Rounding gives
\[
 \E[\ALG]\leq 10\CPB\cdot\OPT
   +(5\CPB+2)k+16,
\]
where $\CPB$ is the absolute constant from Chasing Positive Bodies at resource
augmentation one and covering sparsity two.  The full reduction, rounding
composition, and quantified main theorem are machine-checked in Lean 4 relative
to exact formal interfaces for the two cited source theorems.
Deterministic rational controls and a fresh independent replay accompany the
formal proof.
\end{abstract}

\section{Introduction}

Dynamic Mixture-of-Experts (MoE) inference must decide how many replica GPUs to
assign to each expert as workloads change.  Moving replicas has a switching
cost, while insufficient replication increases bottleneck latency.  Huang,
Lou, and Xiao formalized this tradeoff as an online problem and proved an
\mbox{$O(\sqrt{\log k})$} randomized competitive ratio for the integral primal
problem~\cite{HLX2026}.  Their \mbox{$\Omega(\sqrt{\log k})$} construction concerns an
offset-Fenchel-dual maximization problem.  Weak duality points in the wrong
direction for transferring it to a primal lower bound, and the asymptotic
primal ratio remained open.

Our main result is that the primal problem has a constant randomized
competitive ratio.  The central observation is that the epigraph of
$r/(1+u)$ admits a finite positive-polyhedral inner approximation with constant
distortion.  This permits a reduction to the positive-body chasing theorem of
Bhattacharya, Buchbinder, Levin, and Saranurak~\cite{BBLS2023}.  Two obstacles
must be handled explicitly.  First, the reciprocal epigraph has a continuum of
tangents, whereas the cited theorem takes finite matrices.  Second, charging a
soft service cost to movement suggests resetting an auxiliary height to zero,
but a zero right-hand-side packing row lies outside the normalized source
theorem.  We use a finite tangent grid and geometrically shrinking positive
resets.  Both repairs preserve a horizon-independent constant.

The reduction yields a deterministic fractional path.  Applying the Lazy
Threshold Rounding lemma of Huang, Lou, and Xiao converts it online to an exact
integral allocation with no expected movement loss and at most a factor-three
service loss.

\paragraph{Verification.}
The new reduction is formalized in Lean 4.  The formal theorem treats the two
published ingredients as explicit source interfaces and proves every new MoE
bridge, including the finite tangent construction, event/reset feasibility,
the comparator movement bound, exact-budget projection, service charging,
per-step-to-path rounding telescope, and final constant assembly.
The source interfaces were audited against pinned source archives.  A fresh
independent referee reproduced 24 tests and both canonical artifact hashes and
returned correctness, openness, and major-significance verdicts.

\section{Model and main result}

Fix $m\geq1$ experts and $k\geq1$ replica GPUs.  The integral state space is
\[
 X_{m,k}=\left\{x\in\Z_{\geq0}^{m}:\sum_{i=1}^{m}x_i=k\right\}.
\]
At round $t$, a nonnegative workload vector $r_t\in\R_{\geq0}^{m}$ is
revealed.  The online algorithm chooses $x_t\in X_{m,k}$ and pays
\[
 f_t(x_t)+\lVert x_t-x_{t-1}\rVert_1,
 \qquad
 f_t(x)=\max_{i\in[m]}\frac{r_{t,i}}{1+x_i},
\]
where the integral initial state $x_0$ is part of the instance.  For a fixed
finite workload sequence, $\OPT$ is the minimum total cost over all integral
paths.  We use the fixed-sequence, or oblivious-adversary, convention underlying
the expected-cost rounding result in~\cite{HLX2026}.

For fixed $k$, write $\operatorname{CR}_{\mathrm{rand}}(k)$ for the infimum of
factors $\Gamma$ with the following property: for every $m\geq1$ there is a
constant $A_{m,k}\geq0$ such that, for every integral initial state, some causal
randomized integral algorithm satisfies
$\E[\ALG]\leq\Gamma\cdot\OPT+A_{m,k}$ on every fixed finite workload sequence.
The additive term may depend on $m$ and $k$, but not on the horizon.

Let $\CPB$ denote an absolute competitive constant supplied by Theorem~1.1 of
\cite{BBLS2023} at resource augmentation $\epsilon=1$ and maximum covering-row
support $d=2$.  We absorb into $\CPB$ the constant from reducing a positive body
to positive halfspaces and the standard conversion between upward movement and
full $\ell_1$ movement.

\paragraph{Two properties of the service cost.}
The proof uses two elementary features of $f_t$.  First, its reciprocal terms
are smooth under capacity scaling: for every $r,u\geq0$,
\[
 \frac{r}{1+u/2}\leq2\frac{r}{1+u}.
\]
Thus halving capacity from $2k$ to $k$ while retaining at least half of every
coordinate increases service by at most a factor two.  Second, $f_t$ is the
maximum of univariate coordinate functions.  Its epigraph therefore separates
into constraints involving only the height and one allocation coordinate, and
the tangent discretization preserves covering-row support $d=2$.
Equation~\eqref{eq:cover-row} makes the sparsity explicit, while
\eqref{eq:service-height} applies the capacity-scaling inequality.

\begin{theorem}[Constant upper bound]\label{thm:upper}
For every $m,k\geq1$ and every integral initial state, there is a causal
randomized integral online algorithm such that every fixed finite nonnegative
workload sequence satisfies
\[
 \E[\ALG]
 \leq 10\CPB\cdot\OPT+(5\CPB+2)k+16.
\]
The multiplicative constant is independent of $m$, $k$, and the horizon.
\end{theorem}

Every realized online path is also a feasible offline path.  Scaling an
equal-workload one-round instance therefore makes the additive constant
negligible, so no competitive factor is below one and
$\operatorname{CR}_{\mathrm{rand}}(k)=\Theta(1)$.

\section{Two source theorems}

We use the following exact consequences of published results.

\paragraph{Positive-body chasing.}
A positive body has the form
\[
 K_t=\{z\in\R_{\geq0}^{D}:C^tz\geq\mathbf1,\;P^tz\leq\mathbf1\},
\]
where both matrices are nonnegative.  Theorem~1.1 of~\cite{BBLS2023} gives,
for $\epsilon\in(0,1]$, an
$O(\epsilon^{-1}\log(d/\epsilon))$-competitive online chaser whose covering
constraints remain exact and whose packing constraints may be violated by a
factor $1+\epsilon$; $d$ is the largest covering-row support.  Its movement is
compared with an exact offline path beginning at zero.  Thus, when
$\epsilon=1$ and $d=2$, the movement factor $\CPB$ is absolute, covering rows
remain exact, and packing right-hand sides may be doubled.  Strictly positive
right-hand sides normalize to one without changing row support.  A covering
row with zero coefficients and zero right-hand side is vacuous and may be
deleted.

\paragraph{Lazy Threshold Rounding.}
Lemma~4.4 of~\cite{HLX2026} applies online to any exact-budget fractional MoE
path, using one shared set of random thresholds over all rounds.  Starting from
the given integral initial state, it produces integral exact-budget states
$\bar x_t$ satisfying
\begin{align}
 f_t(\bar x_t)&\leq3f_t(x_t),\label{eq:round-service}\\
 \E\lVert\bar x_t-\bar x_{t-1}\rVert_1
 &\leq\lVert x_t-x_{t-1}\rVert_1.\label{eq:round-move}
\end{align}
The same thresholds are used in consecutive prefixes, so summing
\eqref{eq:round-move} gives the corresponding expected path-movement bound.
The source implements each step greedily in polynomial time.

\section{A finite positive approximation to reciprocal service}

Fix a workload coordinate $r\geq0$ and augmented allocation $u\geq0$.  Write
$q=1+u$.  For a scale $p\geq1$, define the tangent
\[
 L_{r,p}(u)=\frac{r(2p-q)}{p^2}
           =\frac{r(2p-1-u)}{p^2}.
\]

\begin{lemma}[Tangent envelope]\label{lem:tangent}
For all $r\geq0$, $q>0$, and $p>0$,
\[
 \frac rq-L_{r,p}(u)=\frac{r(q-p)^2}{qp^2}\geq0.
\]
If $0\leq u\leq2k$, some $p\in\{1,\ldots,2k+1\}$ also satisfies
\[
 L_{r,p}(u)\geq\frac34\frac r{1+u}.
\]
\end{lemma}

\begin{proof}
The first identity follows by putting the two terms over $qp^2$.  For the
second claim, take $p=\lceil q\rceil$.  Then $1\leq p\leq2k+1$ and
$q\leq p<q+1\leq2q$.  Hence $a=q/p\in(1/2,1]$, and
\[
 \frac{L_{r,p}(u)}{r/q}=a(2-a)\geq\frac34.
\]
The case $r=0$ is immediate.
\end{proof}

Define the finite envelope
\[
 H_r(u)=\max_{1\leq p\leq2k+1}L_{r,p}(u).
\]
Lemma~\ref{lem:tangent} gives
\begin{equation}\label{eq:envelope}
 \frac34\frac r{1+u}\leq H_r(u)\leq\frac r{1+u}
 \qquad(0\leq u\leq2k).
\end{equation}
For $r>0$, the inequality $s\geq L_{r,p}(u)$ is equivalent to the positive
covering row
\begin{equation}\label{eq:cover-row}
 \frac{p^2}{r(2p-1)}s+\frac1{2p-1}u\geq1.
\end{equation}
Its support has size at most two.  This is the $d=2$ sparsity supplied by the
maximum-of-univariate form of $f_t$.  Coordinates with $r=0$ require no row;
equivalently, one may insert and delete a vacuous zero-coefficient covering row
with zero right-hand side, as in the formal reduction.

\begin{remark}
A geometric grid $p\in\{(3/2)^j:(3/2)^j\leq1+2k\}$ gives the same $3/4$
factor with $O(\log(k+1))$ rows per expert.  The concrete Lean reduction uses
the elementary integer grid above; both envelope lemmas are machine-checked.
\end{remark}

\section{The event/reset reduction}

Use augmented coordinates $(u_1,\ldots,u_m,s)\in\R_{\geq0}^{m+1}$.  When
request $r_t$ arrives, present the positive body
\begin{equation}\label{eq:event}
 E_t=\left\{(u,s):\sum_i u_i\leq k,\quad
 s\geq L_{r_{t,i},p}(u_i)\;\text{for all }i,p\right\}.
\end{equation}
After taking the round-$t$ MoE action, present the reset body
\begin{equation}\label{eq:reset}
 R_t=\left\{(u,s):\sum_i u_i\leq k,\quad s\leq\delta_t\right\},
 \qquad \delta_t=2^{-t}\quad(t\geq1),
\end{equation}
and set $\delta_0:=0$ for the initial accounting convention.
Both packing rows have positive right-hand side after normalization.  The event
covering rows have support at most two.  Both bodies are nonempty.  The reset is
processed after the action at time $t$ and before the next request, so the
construction remains causal.

Run the positive-body chaser with $\epsilon=1$.  At an event, write its point
as $(u_t,s_t)$.  At the following reset, write its height as $b_t$.  Then
\begin{equation}\label{eq:aug-feasible}
 \sum_i u_{t,i}\leq2k,\qquad s_t\geq H_{r_{t,i}}(u_{t,i})\quad\forall i,
 \qquad 0\leq b_t\leq2\delta_t.
\end{equation}

\subsection{Removing resource augmentation}

Map every $u\geq0$ with $\sum_i u_i\leq2k$ to the exact-budget fractional
allocation
\begin{equation}\label{eq:balanced}
 \rho(u)=\frac{k-\frac12\sum_i u_i}{m},
 \qquad x_i(u)=\frac{u_i}{2}+\rho(u).
\end{equation}

\begin{lemma}[Balanced projection]\label{lem:balanced}
The map in \eqref{eq:balanced} satisfies $x(u)\geq0$,
$\sum_i x_i(u)=k$, $x_i(u)\geq u_i/2$, and
\[
 \lVert x(u)-x(v)\rVert_1\leq\lVert u-v\rVert_1.
\]
\end{lemma}

\begin{proof}
The budget bound gives $\rho(u)\geq0$, and summing
\eqref{eq:balanced} gives total mass $k$.  Moreover,
\begin{align*}
 \lVert x(u)-x(v)\rVert_1
 &\leq\frac12\lVert u-v\rVert_1+m|\rho(u)-\rho(v)|\\
 &=\frac12\lVert u-v\rVert_1
   +\frac12\left|\sum_i(u_i-v_i)\right|\\
 &\leq\lVert u-v\rVert_1.
\end{align*}
\end{proof}

Combining \eqref{eq:envelope}, \eqref{eq:aug-feasible}, and the capacity-halving
bound $1+x_i(u_t)\geq(1+u_{t,i})/2$ gives
\begin{equation}\label{eq:service-height}
 f_t(x(u_t))
 \leq2\max_i\frac{r_{t,i}}{1+u_{t,i}}
 \leq\frac83s_t.
\end{equation}

\subsection{Charging service to vertical movement}

Let $a_{t-1}$ be the chaser height after the preceding reset, with $a_0=0$.
Thus $a_{t-1}\leq2\delta_{t-1}$ for every $t\geq1$, including the initial
case.  Let $V_t$ be the vertical movement while processing $E_t,R_t$.
Endpoint distances imply
\[
 V_t\geq|s_t-a_{t-1}|+|s_t-b_t|
     \geq2s_t-a_{t-1}-b_t.
\]
Using \eqref{eq:service-height}, $a_{t-1}\leq2\delta_{t-1}$, and
$b_t\leq2\delta_t$, we obtain
\[
 f_t(x(u_t))
 \leq\frac43V_t+\frac83(\delta_{t-1}+\delta_t).
\]
Because $\delta_0=0$ and $\sum_{t\geq1}\delta_t=1$, summing the error terms
contributes at most $16/3$.  If $M$ denotes the chaser's complete movement,
then
\begin{equation}\label{eq:service-bridge}
 \sum_t f_t(x(u_t))\leq\frac43M+\frac{16}{3}.
\end{equation}
The balanced projection is nonexpansive.  Write $x_t:=x(u_t)$ for $t\geq1$
and identify the supplied integral $x_0$ with its real vector.  The first state
can be at distance at most $2k$ from $x_0$, the diameter of the exact simplex.
Hence
\begin{equation}\label{eq:movement-bridge}
 \lVert x_1-x_0\rVert_1+
   \sum_{t=2}^{T}\lVert x_t-x_{t-1}\rVert_1\leq M+2k.
\end{equation}

\section{Comparison with the MoE optimum}

Take any integral comparator path $y_1,\ldots,y_T$ and let
$h_t=f_t(y_t)$.  At event time use the body point $(y_t,h_t)$, and at reset
time use $(y_t,0)$.  The tangent-underestimator identity makes the event point
feasible; the reset point is feasible as well.  Starting from zero, this lifted
body path has movement
\begin{align}
 &k+\sum_{t=2}^{T}\lVert y_t-y_{t-1}\rVert_1+2\sum_{t=1}^{T}h_t
 \notag\\
 &\hspace{35mm}\leq k+2\,\operatorname{cost}(y_1,\ldots,y_T).
 \label{eq:comparator}
\end{align}
Applying the positive-body theorem to an optimal MoE path gives
\begin{equation}\label{eq:M}
 M\leq\CPB(k+2\OPT).
\end{equation}

We now prove the upper bound.

\begin{proof}[Proof of Theorem~\ref{thm:upper}]
Apply Lazy Threshold Rounding to the fractional path constructed above.  Its
two guarantees, together with \eqref{eq:service-bridge} and
\eqref{eq:movement-bridge}, give
\begin{align*}
 \E[\ALG]
 &\leq (M+2k)+3\left(\frac43M+\frac{16}{3}\right)\\
 &=5M+2k+16\\
 &\leq10\CPB\cdot\OPT\\
 &\quad +(5\CPB+2)k+16,
\end{align*}
where the last step uses \eqref{eq:M}.  The additive term is independent of the
horizon, and every step depends only on revealed requests and the one initial
rounding seed.
\end{proof}

\section{Formal and executable verification}

The formal development uses Lean 4.32.2 and Mathlib 4.32.2.  Its main
upper-bound declaration is
\begin{itemize}
\item
\texttt{DynamicMoeMain.dynamicMoe\_explicit\_constant\_upper\_of\_source\_theorems}.
\end{itemize}
It states the explicit factor $10\CPB$.  The proof files also contain the
CPB-compatible body condition and the theorem telescoping HLX's per-step
expected movement inequality.  Lean reports only the standard axioms
\texttt{propext}, \texttt{Classical.choice}, and \texttt{Quot.sound} for every
audited declaration.

The source-relative boundary is explicit: the development does not reprove the
complete CPB or HLX papers.  It formalizes their exact consequences as primitive
premises and proves all new Dynamic MoE implications from them.  Pinned source
archives and theorem-bearing members are checked by SHA-256 before replay.

From the repository root, the complete replay is:
\begin{verbatim}
python3 scripts/bootstrap_dynamic_moe_sources.py
python3 -m pytest tests/test_bootstrap_dynamic_moe_sources.py \
  tests/test_dynamic_moe_positive_body_eval.py \
  tests/test_dynamic_moe_positive_body_formal.py \
  tests/test_dynamic_moe_semantics_formal.py \
  tests/test_dynamic_moe_lower_bound_formal.py \
  tests/test_dynamic_moe_formal_eval.py -q
python3 -m evaluators.dynamic_moe_positive_body_eval \
  --output output/dynamic-moe-serving/positive-body-reduction-control.json
python3 -m evaluators.dynamic_moe_formal_eval \
  --output output/dynamic-moe-serving/formal-reduction-audit.json
\end{verbatim}
The run passes 24 tests.  The canonical artifact hashes are
\begin{align*}
\text{exact control: }&
\texttt{53c7ac33689b4cef22057fd1f9c040dab4e00a8aa13c5277b6a1b6e91e041a66},\\
\text{formal packet: }&
\texttt{adcb13dad31d6fc688982068396e9733cf3d6f6444285f015cd0e6e1c1fde3f4}.
\end{align*}
A fresh independent process reproduced both hashes exactly.  The formal packet
normalizes machine-local path and architecture metadata.  The exact control
includes exhaustive rational tests of the tangent envelope, reset service
charge, exact-budget maps, movement normalization, and a small offline dynamic
program, together with a deliberately corrupted movement negative control.

\section{Scope and open directions}

The theorem resolves the asymptotic randomized integral ratio under the
fixed-sequence convention.  It does not determine the best numerical constant,
the deterministic integral competitive ratio, or the ratio against an adaptive
adversary.  The concrete integer tangent grid has $O(mk)$ covering rows per
event.  The geometric grid noted after Lemma~\ref{lem:tangent} reduces this to
$O(m\log(k+1))$ without changing the proof.  The reset schedule uses exact real
numbers; replacing it with another positive summable schedule can be useful for
an explicit bit-complexity implementation.

\paragraph{Research-system disclosure.}
Beyond, the research system operated by Nth Research Collective, materially
assisted with literature retrieval, hypothesis generation, formalization,
executable controls, and adversarial review.  Ian D'Ambrosio selected the
target, directed the investigation, reconciled the evidence, verified the final
claims, and accepts full responsibility for the manuscript.

\begin{thebibliography}{9}

\bibitem{BBLS2023}
S.~Bhattacharya, N.~Buchbinder, R.~Levin, and T.~Saranurak,
\newblock Chasing positive bodies,
\newblock arXiv:2304.01889v2, 2023 (revised 2024).
\newblock \url{https://arxiv.org/abs/2304.01889v2}.

\bibitem{HLX2026}
Z.~Huang, Q.~Lou, and T.~Xiao,
\newblock Mixture-of-Experts Serving,
\newblock arXiv:2607.17880v1, 2026.
\newblock \url{https://arxiv.org/abs/2607.17880v1}.

\end{thebibliography}

\end{document}
